\documentclass[12pt]{article}
\usepackage{amsmath, amssymb}
\usepackage[margin=1in]{geometry}
\setlength{\parskip}{6pt}
\title{QUBO Formulation: Weighted Max-2-SAT}
\date{}
\begin{document}
\maketitle

\subsection*{Weighted Max-2-SAT}

\paragraph{Problem Statement.}
Given a set of $n$ Boolean variables and a collection $\mathcal{C}$ of weighted clauses, where each clause contains at most two literals, the objective is to assign truth values to all variables so as to maximize the total weight of satisfied clauses. A clause is satisfied if at least one of its literals evaluates to true under the assignment.

\paragraph{Decision Variables.}
Let $x_i \in \{0,1\}$ for $i \in \{1, \dots, n\}$ denote the binary decision variables, where $x_i = 1$ represents the variable being true and $x_i = 0$ represents it being false. Each clause $c \in \mathcal{C}$ has an associated positive weight $w_c > 0$. Each literal in a clause is associated with a variable index and a sign $s \in \{+1, -1\}$, where $+1$ denotes a positive literal ($x_i$) and $-1$ denotes a negated literal ($\neg x_i$).

\paragraph{QUBO Derivation.}
\textbf{Step 1 --- Original Objective.} The problem requires maximizing the total weight of satisfied clauses:
\begin{align}
    \max_{x \in \{0,1\}^n} \sum_{c \in \mathcal{C}} w_c \cdot \mathbb{I}(c \text{ is satisfied}),
\end{align}
where $\mathbb{I}(\cdot)$ is the indicator function. Since QUBO solvers minimize energy, we negate the objective:
\begin{align}
    \min_{x \in \{0,1\}^n} H_0(x) = -\sum_{c \in \mathcal{C}} w_c \cdot \mathbb{I}(c \text{ is satisfied}).
\end{align}

\textbf{Step 2 --- Constraints.} The problem is unconstrained; any binary assignment is feasible. Thus, no constraint terms are introduced.

\textbf{Step 3 --- Algebraic Expansion of Satisfaction.} A clause $c$ is satisfied if at least one literal evaluates to true. Let $L_c$ be the set of literals in clause $c$. For a literal $l \in L_c$ on variable $i$ with sign $s_l$, its truth value is $y_l = \frac{1+s_l}{2}x_i + \frac{1-s_l}{2}(1-x_i)$. The clause is unsatisfied if and only if $\prod_{l \in L_c} (1-y_l) = 1$. Therefore,
\begin{align}
    \mathbb{I}(c \text{ is satisfied}) = 1 - \prod_{l \in L_c} (1-y_l).
\end{align}
Substituting into $H_0$ yields:
\begin{align}
    H_0(x) = \sum_{c \in \mathcal{C}} w_c \prod_{l \in L_c} (1-y_l) - \sum_{c \in \mathcal{C}} w_c.
\end{align}
We expand the product for clauses of size one and two. For a unit clause $c$ with literal $l$ on variable $i$:
\begin{align}
    1-y_l = \frac{1+s_l}{2} - s_l x_i.
\end{align}
For a binary clause $c$ with literals $l_1, l_2$ on variables $i, j$ with signs $s_1, s_2$:
\begin{align}
    (1-y_{l_1})(1-y_{l_2}) = \left(\frac{1+s_1}{2} - s_1 x_i\right)\left(\frac{1+s_2}{2} - s_2 x_j\right) = \frac{(1+s_1)(1+s_2)}{4} - \frac{s_2(1+s_1)}{2}x_i - \frac{s_1(1+s_2)}{2}x_j + s_1 s_2 x_i x_j.
\end{align}

\textbf{Step 4 --- Combined Hamiltonian.} Collecting all linear, quadratic, and constant terms across all clauses, the total energy function $H(x)$ takes the standard QUBO form:
\begin{align}
    H(x) = \sum_{i=1}^n Q_{i,i} x_i + \sum_{1 \le i < j \le n} Q_{i,j} x_i x_j + c.
\end{align}

\textbf{Step 5 --- Q Matrix and Offset.} Reading off the coefficients from the expanded terms gives the general closed-form expressions for the QUBO parameters:
\begin{align}
    Q_{i,i} &= \sum_{c \in \mathcal{C}_1(i)} -w_c s_{l_c} + \sum_{j \neq i} \sum_{c \in \mathcal{C}_2(i,j)} -\frac{s_{c,j}(1+s_{c,i})}{2} w_c, \\
    Q_{i,j} &= \sum_{c \in \mathcal{C}_2(i,j)} w_c s_{c,i} s_{c,j} \quad \text{for } i < j, \\
    c &= -\sum_{c \in \mathcal{C}} w_c + \sum_{c \in \mathcal{C}_1} w_c \frac{1+s_{l_c}}{2} + \sum_{c \in \mathcal{C}_2} w_c \frac{(1+s_{c,i})(1+s_{c,j})}{4},
\end{align}
where $\mathcal{C}_1(i)$ denotes unit clauses containing variable $i$, $\mathcal{C}_2(i,j)$ denotes binary clauses containing both variables $i$ and $j$, and $s_{c,i}, s_{c,j}$ are the signs of the literals on variables $i$ and $j$ in clause $c$.

\paragraph{Penalty Weight Justification.}
The Weighted Max-2-SAT problem is inherently unconstrained; every binary assignment corresponds to a valid truth valuation. Consequently, no penalty terms or penalty weights are required. The QUBO formulation directly encodes the objective function without auxiliary variables or constraint penalties.

\paragraph{Correctness Argument.}
Minimizing $H(x)$ over $\{0,1\}^n$ recovers the optimal assignment because the energy function is constructed as a linear transformation of the clause-satisfaction weight. For any assignment $x$, the term $\sum_{c \in \mathcal{C}} w_c \prod_{l \in L_c} (1-y_l)$ evaluates exactly to the total weight of unsatisfied clauses. Thus, $H(x) = W_{\text{unsat}}(x) - W_{\text{total}}$, where $W_{\text{total}} = \sum_{c \in \mathcal{C}} w_c$ is constant. Minimizing $H(x)$ is therefore equivalent to minimizing the weight of unsatisfied clauses, which is mathematically identical to maximizing the weight of satisfied clauses. Since the domain is unconstrained, the global minimum of $H(x)$ must correspond to a feasible assignment that achieves the maximum possible satisfied weight.

\paragraph{Illustrative Example.}
Consider an instance with $n=3$ variables and three clauses:
\begin{itemize}
    \item $c_1: (x_1 \lor \neg x_2)$ with weight $w_1 = 3$. Literals have signs $s_{1,1}=+1, s_{1,2}=-1$.
    \item $c_2: (\neg x_1 \lor x_3)$ with weight $w_2 = 2$. Literals have signs $s_{2,1}=-1, s_{2,3}=+1$.
    \item $c_3: (x_2)$ with weight $w_3 = 4$. Literal has sign $s_{3,2}=+1$.
\end{itemize}
Applying the general formulas:
\begin{align}
    Q_{1,1} &= -\frac{(-1)(1+1)}{2}(3) - \frac{(1)(1-1)}{2}(2) = 3, \\
    Q_{2,2} &= -(4)(1) - \frac{(1)(1-1)}{2}(3) = -4, \\
    Q_{3,3} &= -\frac{(-1)(1+1)}{2}(2) = 2, \\
    Q_{1,2} &= (3)(+1)(-1) = -3, \quad Q_{1,3} = (2)(-1)(+1) = -2, \quad Q_{2,3} = 0.
\end{align}
The constant offset is:
\begin{align}
    c = -(3+2+4) + (4)\frac{1+1}{2} + (3)\frac{(2)(0)}{4} + (2)\frac{(0)(2)}{4} = -9 + 4 = -5.
\end{align}
The resulting QUBO energy function is:
\begin{align}
    H(x) = 3x_1 - 4x_2 + 2x_3 - 3x_1x_2 - 2x_1x_3 - 5.
\end{align}
Evaluating $H(x)$ over all $2^3=8$ assignments yields the minimum energy $H(1,0,1) = -5$. This corresponds to the assignment $x_1=1, x_2=0, x_3=1$, which satisfies $c_1$ and $c_2$ (total weight $3+2=5$) and leaves $c_3$ unsatisfied, confirming the formulation correctly recovers the optimal solution.

\end{document}